\begin{frame}{같은 구조, 다른 질문 — 한 학기가 남기는 선물}
  \begin{block}{4단계는 \textbf{그대로} 살아 있다}
    ① 모델 정의 $\to$ ② 손실 $\to$ ③ $\theta$ 조정 $\to$ ④
    통제는 ML2에서도 \textbf{그대로} — 바뀌는 것은 \textbf{
    ①의 ``모델''}(\textbf{정책 $\pi(a|s)$})과 \textbf{④의
    ``데이터''}(\textbf{에이전트의 경험})\textbf{뿐}.
  \end{block}
  \vskip0.4em
  \kq{\textbf{같은 구조, 다른 질문} — 이번 학기가 ML2에 남기는
  \textbf{정확한 선물}. \; 어떤 모델을 쓸지보다 \textbf{먼저}
  ``\textbf{무엇을 데이터로 갖고, 무엇을 맞히고 싶은가}.''}
  \vskip0.5em
  \small
  \begin{itemize}
    \item \textbf{Q. 기말고사 대비 — 어디서부터?} \; 개념 연결
      4문항 + 절차 3문항을 그 순서로. \textbf{6번(누수 질문)}을
      노트북 없이 답하면 준비 끝.
    \item \textbf{Q. Ch02~15가 섞여 있는데 기말은 Block B
      만?} \; 복습 문제가 \textbf{의도적으로} 두 블록을 넘나
      듦 — ``\textbf{연결을 아는가}''가 ``개별 장을 외우는 것
      ''보다 시험·ML2에서 훨씬 많이 쓰임.
    \item \textbf{Q. ML1의 수식 유도를 전부 외워야?} \;
      \textbf{아니오} — 4단계 구조의 \textbf{감각}이 자동인 것
      이 목표. \textbf{유도}(GDA$\leftrightarrow$로지스틱,
      ELBO)는 필요할 때 장에서 \textbf{찾아보는 것}이 정상,
      ``\textbf{손실을 뭐로 나타내지?}''가 \textbf{자동 반사}
      여야.
    \item \textbf{Q. ML2의 프로젝트는 뭐가 다른가요?} \; 뼈
      대(발표+리뷰+채점)는 \textbf{동일}. ``\textbf{결과의
      정직성}''의 \textbf{초점}만 바뀜 — RL에서 ``학습이 안
      된 부분''은 실패가 아니라 \textbf{데이터} — 왜 안 됐
      는지 Ch09~15의 개념으로 진단하는 것이 \textbf{좋은
      보고서}.
  \end{itemize}
\end{frame}
